% GENERATED FILE. Edit canonical lessons, then run npm run generate:books.
% canonical-source-hash: fnv1a64:e09c59248038c676
% canonical-lessons: IT-C425-terminal, IT-C425-scalo, IT-C425-imbarco, IT-C425-navetta, IT-C425-fuso-orario

\chapter{In aeroporto}
\label{ch:it-in-aeroporto}
\hlchaptermodality{425}

\begin{chapteropening}
\textbf{By the end of this chapter:} \emph{I can find my way through an airport: the terminal, a stopover, boarding, the shuttle bus and the time zone.}

Complete the last lesson of chapter 425: I can find my way through an airport: the terminal, a stopover, boarding, the shuttle bus and the time zone.
\end{chapteropening}

\section[il terminal]{il terminal — the terminal (at an airport)}

\label{lesson:IT-C425-terminal}



\begin{quote}
Before the new one: say the Italian for the lifebelt, then the Italian for the nausea.
\end{quote}

\subsection*{You'll want to know: il terminal}
\textbf{il terminal} — \textquotedblleft{}the terminal (at an airport)\textquotedblright{}.

\begin{quote}
Il nostro volo parte dal terminal due.
\end{quote}

\begin{quote}
\emph{Our flight leaves from terminal two.}
\end{quote}

Taken from English; masculine, unchanged in the plural.

\begin{cousinweb}
\emph{terminal} was borrowed from English \emph{terminal}, which goes back to Latin \emph{terminālis} 'of the end', from \emph{terminus} 'boundary, end' (EtymDB).
\end{cousinweb}

\begin{grammarlens}[title={What you've built}]
One more word for telling, explaining and describing at B1.
\end{grammarlens}

\subsection*{Your turn}
\begin{itemize}
\raggedright
  \item \emph{Say it:} \emph{il terminal}
  \item \emph{Say it:} \emph{il terminal}, once more
  \item \emph{Say it:} a sentence of your own with \emph{il terminal}
  \item \emph{Recall:} say the Italian for the lifebelt, then the Italian for the nausea, then say \emph{il terminal} again
\end{itemize}

\begin{tcolorbox}[breakable,colback=teal!4,colframe=teal!35!black,title={Before you move on}]
What does il terminal mean? (\textquotedblleft{}The terminal (at an airport)\textquotedblright{}.) Say it once more.
\end{tcolorbox}
\section[lo scalo]{lo scalo — the stopover; the port of call}

\label{lesson:IT-C425-scalo}



\begin{quote}
Before the new one: say the Italian for the nausea, then the Italian for the terminal.
\end{quote}

\subsection*{You'll want to know: lo scalo}
\textbf{lo scalo} — \textquotedblleft{}the stopover; the port of call\textquotedblright{}.

\begin{quote}
Il volo fa uno scalo di due ore.
\end{quote}

\begin{quote}
\emph{The flight has a two-hour stopover.}
\end{quote}

\emph{Un volo senza scalo} — a direct flight; \emph{fare scalo} — to stop over.

\begin{cousinweb}
\emph{scalo} is built on \emph{scala} 'ladder, steps', which is inherited from Latin \emph{scāla} (EtymDB): first the steps where ships landed goods.
\end{cousinweb}

\begin{grammarlens}[title={What you've built}]
One more word for telling, explaining and describing at B1.
\end{grammarlens}

\subsection*{Your turn}
\begin{itemize}
\raggedright
  \item \emph{Say it:} \emph{lo scalo}
  \item \emph{Say it:} \emph{lo scalo}, once more
  \item \emph{Say it:} a sentence of your own with \emph{lo scalo}
  \item \emph{Recall:} say the Italian for the nausea, then the Italian for the terminal, then say \emph{lo scalo} again
\end{itemize}

\begin{tcolorbox}[breakable,colback=teal!4,colframe=teal!35!black,title={Before you move on}]
What does lo scalo mean? (\textquotedblleft{}The stopover; the port of call\textquotedblright{}.) Say it once more.
\end{tcolorbox}
\section[l'imbarco]{l'imbarco — the boarding}

\label{lesson:IT-C425-imbarco}



\begin{quote}
Before the new one: say the Italian for the terminal, then the Italian for the stopover.
\end{quote}

\subsection*{You'll want to know: l'imbarco}
\textbf{l'imbarco} — \textquotedblleft{}the boarding\textquotedblright{}.

\begin{quote}
L'imbarco chiude venti minuti prima.
\end{quote}

\begin{quote}
\emph{Boarding closes twenty minutes before.}
\end{quote}

\emph{La carta d'imbarco} — the boarding pass.

\begin{cousinweb}
\emph{imbarco} is built on \emph{imbarcare} 'to take on board', from \emph{in-} and \emph{barca} 'boat', which is inherited from Late Latin \emph{barca}.
\end{cousinweb}

\begin{grammarlens}[title={What you've built}]
One more word for telling, explaining and describing at B1.
\end{grammarlens}

\subsection*{Your turn}
\begin{itemize}
\raggedright
  \item \emph{Say it:} \emph{l'imbarco}
  \item \emph{Say it:} \emph{l'imbarco}, once more
  \item \emph{Say it:} a sentence of your own with \emph{l'imbarco}
  \item \emph{Recall:} say the Italian for the terminal, then the Italian for the stopover, then say \emph{l'imbarco} again
\end{itemize}

\begin{tcolorbox}[breakable,colback=teal!4,colframe=teal!35!black,title={Before you move on}]
What does l'imbarco mean? (\textquotedblleft{}The boarding\textquotedblright{}.) Say it once more.
\end{tcolorbox}
\section[la navetta]{la navetta — the shuttle (bus)}

\label{lesson:IT-C425-navetta}



\begin{quote}
Before the new one: say the Italian for the stopover, then the Italian for the boarding.
\end{quote}

\subsection*{You'll want to know: la navetta}
\textbf{la navetta} — \textquotedblleft{}the shuttle (bus)\textquotedblright{}.

\begin{quote}
Una navetta va dall'aeroporto all'albergo.
\end{quote}

\begin{quote}
\emph{A shuttle goes from the airport to the hotel.}
\end{quote}

Literally a little ship; also the shuttle of a loom.

\begin{cousinweb}
\emph{navetta} is a little \emph{nave} 'ship', which is inherited from Latin \emph{nāvis} (Lewis \& Short).
\end{cousinweb}

\begin{grammarlens}[title={What you've built}]
One more word for telling, explaining and describing at B1.
\end{grammarlens}

\subsection*{Your turn}
\begin{itemize}
\raggedright
  \item \emph{Say it:} \emph{la navetta}
  \item \emph{Say it:} \emph{la navetta}, once more
  \item \emph{Say it:} a sentence of your own with \emph{la navetta}
  \item \emph{Recall:} say the Italian for the stopover, then the Italian for the boarding, then say \emph{la navetta} again
\end{itemize}

\begin{tcolorbox}[breakable,colback=teal!4,colframe=teal!35!black,title={Before you move on}]
What does la navetta mean? (\textquotedblleft{}The shuttle (bus)\textquotedblright{}.) Say it once more.
\end{tcolorbox}
\section[il fuso orario]{il fuso orario — the time zone}

\label{lesson:IT-C425-fuso-orario}



\begin{quote}
Before the new one: say the Italian for the boarding, then the Italian for the shuttle.
\end{quote}

\subsection*{You'll want to know: il fuso orario}
\textbf{il fuso orario} — \textquotedblleft{}the time zone\textquotedblright{}.

\begin{quote}
Tra qui e casa ci sono sei ore di fuso orario.
\end{quote}

\begin{quote}
\emph{There are six hours of time difference between here and home.}
\end{quote}

Plural \emph{i fusi orari}. A \emph{fuso} is a spindle: on a globe the zones look like spindles.

\begin{cousinweb}
\emph{fuso orario} is built from \emph{fuso} 'spindle', inherited from Latin \emph{fūsus} (EtymDB), and \emph{orario} 'of the hours', from Latin \emph{hōrārius}, built on \emph{hōra} 'hour'.
\end{cousinweb}

\begin{grammarlens}[title={What you've built}]
One more word for telling, explaining and describing at B1.
\end{grammarlens}

\subsection*{Your turn}
\begin{itemize}
\raggedright
  \item \emph{Say it:} \emph{il fuso orario}
  \item \emph{Say it:} \emph{il fuso orario}, once more
  \item \emph{Say it:} a sentence of your own with \emph{il fuso orario}
  \item \emph{Recall:} say the Italian for the boarding, then the Italian for the shuttle, then say \emph{il fuso orario} again
\end{itemize}

\begin{tcolorbox}[breakable,colback=teal!4,colframe=teal!35!black,title={Before you move on}]
What does il fuso orario mean? (\textquotedblleft{}The time zone\textquotedblright{}.) Say it once more.
\end{tcolorbox}
